\section{Introduction}
%The original Transformer structure was designed for machine translation that uses an \textit{Encoder} to represent the source sentence and a \textit{Decoder} to generate a target sentence.  The \textit{Encoder} and \textit{Decoder} structure in Transformer could be leveraged separately for different purpose.  Almost all recent breakthroughs in language model pre-training benefit from the Transformer structure, and they could be categorized into three classes based on how the transformer structures are used as the building blocks. Some \citepp{devlin2018bert} only uses the \textit{Encoder} to build an auto-encoder language model, remarkably the masked language model in BERT.  Some \citepp{gpt22019} uses its decoder to pre-train an auto-regressive language model, such as GPT-2. The other formalize the language modeling as a seq-to-seq generation task and use both the \textit{Encoder} and \textit{Decoder} as the original Transformer.   
%One of the design choices in Transformer is the modeling of the position information in a sequence, since its self-attention mechanism is insensitive to the order of the content in the sequence. Therefore, a Transformer model usually adapts a positional bias to mimic the sequential dependency. There are usually two kinds of positional bias. One is relative positional bias as in the original transformer paper \citep{vaswani2017attention}. The other is absolute positional bias as GPT \citep{gpt22019}, BERT \citep{devlin2018bert} and RoBERTa \citep{liu2019roberta}.  Both of them perform an element-wise summation to mix the positional bias with the content embedding directly. We argue that this entangle design will inevitably introduce noise information into the input data and hurt the performance. 
%However, as the positional bias is independent of content, adding position bias directly into the content will introduce some noise to the content and it will also weaken the position dependencies of tokens in the sequence.
The Transformer has become the most effective neural network architecture for neural language modeling. 
Unlike recurrent neural networks (RNNs) that process text in sequence, Transformers apply self-attention to compute in parallel every word from the input text an attention weight that gauges the influence each word has on another, thus allowing for much more parallelization than RNNs for large-scale model training \citep{vaswani2017attention}. 
Since 2018, we have seen the rise of a set of large-scale Transformer-based Pre-trained Language Models (PLMs), such as 
GPT \citep{radford2019language,brown2020language},
%GPT-2 \citep{radford2019gpt2},
%GPT-3 \citep{brown2020language}, 
BERT \citep{devlin2018bert}, 
RoBERTa \citep{liu2019roberta}, 
XLNet \citep{yang2019xlnet}, 
UniLM \citep{dong2019unilm},
ELECTRA \citep{clark2020electra},
T5 \citep{raffel2019t5},
ALUM \citep{liu2020alum},
StructBERT \citep{wang2019structbert} and
ERINE \citep{sun2019ernie}
. 
These PLMs have been fine-tuned using task-specific labels and created new state of the art in
many downstream natural language processing (NLP) tasks \citep{liu2019mt-dnn,minaee2020deep, jiang2019smart,he2019hnn,he2019x,shen2020exploiting}. 

In this paper, we propose a new Transformer-based neural language model \textbf{DeBERTa} (\textbf{D}ecoding-\textbf{e}nhanced \textbf{BERT} with disentangled \textbf{a}ttention), which 
improves previous state-of-the-art PLMs using two novel techniques: a disentangled attention mechanism, and an enhanced mask decoder.
%, and a virtual adversarial training method for fine-tuning.
%which has proven more effective than RoBERTa and BERT and after fine-tuning leads to better results on a wide range of NLP tasks. 
%DeBERTa makes two modifications to the BERT model.

\paragraph{Disentangled attention.}
% DeBERTa uses a disentangled self-attention mechanism. 
Unlike BERT where each word in the input layer is represented using a vector which is the sum of its word (content) embedding and position embedding, each word in DeBERTa is represented using two vectors that encode its content and position, respectively, and the attention weights among words are computed using disentangled matrices based on their contents and relative positions, respectively. This is motivated by the observation that the attention weight of a word pair depends on not only their contents but their relative positions. For example, the dependency between the words ``deep'' and ``learning'' is much stronger when they occur next to each other than when they occur in different sentences. 

\paragraph{Enhanced mask decoder.}
Like BERT, DeBERTa is pre-trained using masked language modeling (MLM). MLM is a fill-in-the-blank task, where a model is taught to use the words surrounding a mask token to predict what the masked word should be. DeBERTa uses the content and position information of the context words for MLM. The disentangled attention mechanism already considers the contents and relative positions of the context words, but not the absolute positions of these words, which in many cases are crucial for the prediction.
Consider the sentence “a new store opened beside the new mall” with the italicized words “store” and “mall” masked for prediction. Although the local contexts of the two words are similar, they play different syntactic roles in the sentence. (Here, the subject of the sentence is “store” not “mall,” for example.) These syntactical nuances depend, to a large degree, upon the words’ absolute positions in the sentence, and so it is important to account for a word’s absolute position in the language modeling process. DeBERTa incorporates absolute word position embeddings right before the softmax layer where the model decodes the masked words based on the aggregated contextual embeddings of word contents and positions.

%As an extension to disentangled attention, we enhance the output layer of BERT for pre-training to address a limitation of relative positions.  We observe in some situations, it is challenging for the relative positions only mechanism to accurately predict masking tokens. 
%For example, 
%Consider a sentence ``\textit{A new \textbf{store} opened near the new \textbf{mall}}'' with the words \textit{store} and \textit{mall} masked for prediction.
%Although the local contexts of the two words are similar, they play different syntactic roles in the sentence depending to a large degree upon their absolute positions in the sentence. 
%DeBERTa incorporates absolute word position embeddings in the softmax layer where the model decodes the masked words based on the aggregated contextual embeddings of word contents and positions.

%only using the relative positions is not sufficient for the model to accurately predict \textit{store} and \textit{mall} in this sentence, since they are exchangeable in syntax although with a different meaning, also right after the word \textit{new} with an exact relative position to it. To address this limitation, we propose to introduce absolute positions back in the output layer of BERT, as a complement to the relative positions.

%The output softmax layer of BERT is replaced with an Enhanced Mask Decoder (\DecoderName) to predict the masked tokens during model pre-training. 
%This is motivated by mitigating a mismatch between pre-training and fine-tuning. While fine-tuning, we use a task-specific decoder that takes the BERT output as input and produces the task labels. 
%However, while pre-training we do not use any task-specific decoder and instead simply normalize the BERT output (logits) via softmax. 
%We thus treat the pre-training task, masked language model (MLM), the same as any fine-tuning task, and add a task-specific decoder, which is implemented as a two-layer Transformer decoder and a softmax output layer, for pre-training.

% In addition, we also make a small but important modification on BERT's output vectors before we feed them into the EMD for prediction. 
%The authors of BERT \citep{devlin2018bert} propose to not replace all masked tokens with the \texttt{[MASK]} token, but keep 10\% of them unchanged. 
%Although this is motivated by mitigating the mismatch between fine-tuning and pre-training since \texttt{[MASK]} never occurs in the input of downstream tasks, the method suffers from information leaking i.e., predicting a token conditioned on the token itself. To address the issue, we replace the output vectors of the 10\% masked but unchanged tokens with their position embedding vectors, inspired by the use of self-exlcusive pattern in XLNet \citep{yang2019xlnet}. 

% BERT proposed to keep a portion of masked tokens, e.g. 10\%, unchanged in the masked language model (MLM) objective to mitigate the mismatch between pre-training and fine-tuning. However, this will encourage the self-attention to look at itself. Self-exclusion pattern was found in XLNet\citep{yang2019xlnet} which outperforms BERT on a wide range of downstream NLU tasks. We argue that to make the model efficient the self-attention should pay more attention to other tokens instead of itself to capture global contextual information. Inspired by two-stream attention in XLNet \citep{yang2019xlnet}, we propose \DecoderName \ which only uses the position and other tokens to decode unchanged mask tokens but without additional parameters and computation cost.  We will show in our experiment, this can effectively encourage the self-attention to capture global contextual information and improve the fine-tuning performance of downstream tasks.

%This is motivated by mitigating a mismatch between pre-training and fine-tuning. In fine-tuning, we often add a task-specific decoder that takes the BERT output as input and produces the task labels. However, in pre-training, we do not use any task-specific decoder but simply normalize the BERT output (logits) via softmax. We thus treat the pre-training task, masked language model (MLM), the same way as any fine-tuning tasks, and add a task-specific decoder, which is implemented as a two-layer Transformer classifier, for pre-training. 

In addition, we propose a new virtual adversarial training method for fine-tuning PLMs to downstream NLP tasks. The method is effective in improving models’ generalization.

%\paragraph{Scale invariant fine-tuning.} 
%Virtual adversarial training is a regularization method for improving models’ generalization~\citep{miyato2018vat,jiang2019smart}. 
%It does so by improving a model’s robustness to adversarial examples, which are created by making small perturbations to the input. The model is regularized so that when given a task-specific example, the model produces the same output distribution as it produces on an adversarial perturbation of that example. For NLP tasks, the perturbation is applied to the word embedding instead of the original word sequence. However, the value ranges (norms) of the embedding vectors vary among different words and models. The variance gets larger for bigger models with billions of parameters, leading to some instability of adversarial training. Inspired by layer normalization~\citep{ba2016layer}, to improve the training stability, we developed a Scale-Invariant-Fine-Tuning (SiFT) method where the perturbations are applied to the \emph{normalized} word embeddings.

We show through a comprehensive empirical study that these techniques substantially improve the efficiency of pre-training and the performance of downstream tasks. 
In the NLU tasks, compared to RoBERTa-Large, a DeBERTa model trained on half the training data performs consistently better on a wide range of NLP tasks, 
achieving improvements on MNLI by +0.9\% (90.2\% vs. 91.1\%), on SQuAD v2.0 by +2.3\%(88.4\% vs. 90.7\%), and RACE by +3.6\% (83.2\% vs. 86.8\%). 
In the NLG tasks, DeBERTa reduces the perplexity from 21.6 to 19.5 on the Wikitext-103 dataset. 
We further scale up DeBERTa by pre-training a larger model that consists of 48 Transformer layers with 1.5 billion parameters. 
The single 1.5B-parameter DeBERTa model substantially outperforms T5 with 11 billion parameters on the SuperGLUE benchmark \citep{wang2019superglue} by 0.6\%(89.3\% vs. 89.9\%), and surpasses the human baseline (89.9 vs. 89.8) for the first time.
The ensemble DeBERTa model sits atop the SuperGLUE leaderboard as of January 6, 2021, outperforming the human baseline by a decent margin (90.3 versus 89.8). 
 

%Transformer has become a full-fledged model structure for natural language pre-training and spawned many breakthroughs in natural language processing (NLP) tasks.  Recent state-of-the-art models engage in improving the vanilla Transformer structure to make it more capable, efficient and effective. For example, Transformer-XL \citep{dai2019transformer} and Music-Transformer \citep{huang2018musictf} improved the original Transformer to tackle long sequence.  ALBERT \citep{lan2019albert} improved Transformer via layer-wise parameter sharing to reduce overall model size. The extensive adaptation of Transformer has resulted in a downsizing of inventing new model structures, but an upsizing of various improvements upon the Transformer structure. In line with previous works, this paper puts forward two improvements to enhance existing Transformer structure.
%achieve a better training efficiency of a Transformer model\xiaodl{Also better performance?}.  
%The first is a more thorough capture of token dependency in a sequence; the other is a novel approach to reconstruct masked tokens in the Masked Language Model (MLM) objective. 

%Given a sequence with multiple tokens, the self-attention mechanism in Transformer lacks of a natural way to seize its sequential information.  Existing approaches usually adapt a positional bias to mimic this sequential dependency. Two kinds of remedies have been extensively applied. One is the absolute positional bias as in the original Transformer paper \citepp{vaswani2017attention}, GPT \citepp{gpt22019}, BERT \citepp{devlin2018bert}, and RoBERTa \citepp{liu2019roberta}, which perform an element-wise summation to blend the positional embedding with the content embedding. The other is the relative positional bias \citepp{huang2018musictf,yang2019xlnet}, which tried to learn separate positional parameters to capture relative positions, partially ignoring the complicated interactions between content and positions. We argue that either this entangled simplification or independent positional parameterizations will inevitably introduce miscellaneous uncertainty into the input data and can result in a sheer performance loss. To address this limitation, we propose a novel approach to efficiently disentangle the positional information from the content information via parameterizing them independently. This produces four different types of cross attentions between content and position, resulting in a more thorough characterization of the interactions between content and positions for each layer in a stacked Transformer. 

%Although original Transformer is designed for seq-to-seq (s2s) tasks with a $Encoder$+$Decoder$ structure for text generation, recent works in Natural Language Understanding (NLU) often leverage its $Encoder$ side for pre-training a language model with a self-supervision objective, such as the Masked Language Model (MLM) \citepp{devlin2018bert, liu2019roberta, lan2019albert}. In this paper, we are still focused on the MLM objective for NLU tasks. However, we adapt a different view on a $N$-layer stacked $Encoder$ Transformer, denoted as $Transformer(Q,K,V)$ as in \citepp{vaswani2017attention}. We first treat the lower $N-1$ layers as the $\mbox{E-Encoder}$ while the top layer as a single-step $\mbox{E-Decoder}$. This $\mbox{E-Decoder}$ is used to produce token-wise contextual embeddings which are used to reconstruct the masked tokens in MLM, rather than traditional auto-regressive text generation. We then extend the $\mbox{E-Decoder}$ from single-step to multi-step via an attention mechanism. This is also inspired by the $encoder-decoder$ attention in the original Transformer model, in which the encoder side provides a static $K$ and $V$, while the $Decoder$ updates the $Q$ when it moves up to next decoder layer. Similarly, during the multi-step decoding,  the $K$ and $V$ will be always taken from the last layer in $\mbox{E-Encoder}$ while we will update the $Q$ in each $\mbox{E-Decoder}$ step. Existing $Q$ will be used  for either generating the new $Q$ in next step, or reconstructing the masked tokens in the last decoding step. By using a dynamic $Q$ to recursively interact with static $K$ and $V$ from the last layer in $\mbox{E-Encoder}$, we conjecture that there are two advantages to achieve a better training efficiency. First, compared with the single-step approach, the final $Q$ can have a deeper understanding of the original $K$ and $V$ from multiple different perspectives, similar to the idea of multi-step reasoning \citepp{shen2017reasonet,liu2018san}. This can lead to better predict the masked tokens. Second, it can push more objective-oriented information back to the static $K$ and $V$ during the back-propagation, since they have interacted with $Q$ multiple times in the forward propagation, making their accumulating gradients better capture the feedback signal in the objective function from all the steps. This can help to learn a better representation for all the $N-1$ layers in $\mbox{E-Encoder}$. Moreover, we design this multi-step decoding methodology without introducing any new parameters into the model or eligible extra computation overhead.  \xiaodl{Need to explain that after pre-training, should we use N-1 layers or N layers.}

%\xiaodl{A mini discussion on pre-training LM?}
%---------------------------------------

%EXPERIMENTAL RESULT COME NEXT

%Paragraph 3: start with the encoder-decoder structure with teacher-forcing, and the purpose is to learn a better encoder representation. 
